%------------------------------------------------------------------------------%
\begin{question}
  Considerando a equação
  \[
    x^2 + 4y^2 - 4x + 8y + 8 = 0
  \]
  \vspace{-\baselineskip}
  \begin{enumerate}
    \item identifique a curva
    \item determine seus principais valores e pontos
    \item esboce a curva
  \end{enumerate}
\end{question}

%------------------------------------------------------------------------------%
{
\definecolor{cx}{RGB}{  0, 0, 255}
\definecolor{cy}{RGB}{255, 0, 255}
\begin{resolution}{Completando Quadrados}
  \begin{align*}
    \onslide<+->{x^2 + 4y^2 - 4x + 8y + 8 & = 0}
    \\
    \onslide<+->{  \mathcolor{cx}{ x^2 - 4x} \hphantom{+4}\;\,}
    \onslide<+->{+ \mathcolor{cy}{4y^2 + 8y} \hphantom{+9}\quad\;}
    \onslide<+->{& = -8}
    \\
    \onslide<+->{  \mathcolor{cx}{  x^2 -  4x} \hphantom{+4}\;\,
                 + \mathcolor{cy}{4(y^2 + 2y)} \hphantom{+9}\;\, & = -8}
    \\
    \onslide<+->{  \mathcolor{cx}{  x^2 - 4x + 4}}
    \onslide<+->{+ \mathcolor{cy}{4(y^2 + 2y + 1)}}
    \onslide<6->{& = -8}
    \onslide<6->{\mathcolor{cx}{+ 4 }}
    \onslide<7->{\mathcolor{cy}{+ 4 \times 1 }}
    \\
    \onslide<+->{\mathcolor{cx}{ (x-2)^2} \hspace{3mm} + \hspace{4mm}
                 \mathcolor{cy}{4(y+1)^2} \hspace{5mm} & = 0}
  \end{align*}
\end{resolution}
}

%------------------------------------------------------------------------------%
\begin{resolution}{Identificação}
  Temos
  \[
    (x-2)^2 + 4(y+1)^2 = 0
  \]
  \pause
  Como
  \[
    (x-2)^2 \geq 0
    \qquad\qquad
    (y+1)^2 \geq 0
  \]
  \pause
  a única solução é
  \[
    x = 2
    \qquad\qquad
    y = -1
  \]

  \pause
  Cônica degenerada é constituída por um único ponto
  \[
    P = (2, -1)
  \]
\end{resolution}

%------------------------------------------------------------------------------%
\begin{resolution}{Esboço}
  \centering
  \includegraphics{I-07-Identificacao-ex-C}
\end{resolution}

%------------------------------------------------------------------------------%
